\appendix
\section{Appendix: code map, report shape, and references}
\label{part:appendix}

\subsection{Where things live}

All paths below are relative to \fx{prototypes/compiler\_harness/} in the
repository root.

\begin{center}
\begin{tabular}{@{}p{0.30\linewidth} p{0.60\linewidth}@{}}
\toprule
Path & Contents \\
\midrule
\fx{fixtures/}
  & 54 case directories, each pairing one MLIR shape with its snapshot; family
    directories also own C evidence and optional case-local candidate bundles \\
\fx{c/}
  & shared build/check support and the executable equivalence driver \\
\fx{c/check\_harness.py}
  & structural validation for snapshots, evidence paths, and case-local
    candidates \\
\fx{integration/}
  & metadata, witness-driver, import, and code-generation tests \\
\fx{integration/candidate-bundles/}
  & reproducibility and negative tests for the nine nonclaimable candidates \\
\fx{examples/actors/}
  & actor and ACL design examples; outside the enforced corpus \\
\fx{FIXTURE\_REVIEW\_GUIDE.md}
  & review order, fixture intent, corpus counts, and manual sign-off criteria \\
\bottomrule
\end{tabular}
\end{center}

Two orientation notes for a reader opening the corpus for the first time. A
case's sibling \fx{snapshot.yaml} is the authoritative preflight boundary and
expectation; its \fx{c\_evidence} paths record provenance, not a claim that the
MLIR was compiled from that C. The fixture inventory is \emph{closed}: every
MLIR case has exactly one sibling snapshot and every snapshot has exactly one
sibling MLIR file, while every family C source is cited by at least one case.

\subsection{Running it}

\begin{center}
\begin{tabular}{@{}p{0.42\linewidth} p{0.5\linewidth}@{}}
\toprule
Command & Effect \\
\midrule
\fx{make check}            & all IR and integration tests \\
\fx{make check-shape}      & all 54 fixture MLIR/snapshot shape tests \\
\fx{make check-mlir}       & compatibility alias for \fx{check-shape} \\
\fx{make check-candidates} & the nine candidate-bundle tests \\
\fx{make check-artifacts}  & compatibility alias for \fx{check-candidates} \\
\fx{make check-integration}& metadata, witness, import, code generation \\
\fx{make bootstrap-lit}    & create the pinned local test runner \\
\bottomrule
\end{tabular}
\end{center}

A clean default suite is the correct outcome: executable preflight checks pass,
while future SPS semantic tests remain unsupported until a Rev-4 verifier and
canonical materialized bundle are supplied. Neither a green preflight test nor
an unsupported semantic arm computes \fx{ModelStatus}.

\subsection{The shape of a report}

The walkthrough implies a record with more structure than a single verdict token.
Collecting what each step required:

\begin{center}
\begin{tabular}{@{}l p{0.60\linewidth}@{}}
\toprule
Field & Constraint that forces it \\
\midrule
model status
  & \Proved{} carries no reason; a counterexample carries a witness; a refusal
    carries an ordered blocker set (\cref{step:verdict}) \\[3pt]
deployment status
  & an independent axis: a model counterexample cannot be repaired by deployment
    evidence (\cref{step:deployment}) \\[3pt]
policy review status
  & a third axis that by construction never changes either of the other two \\[3pt]
diagnostic disposition
  & per site, from the five-valued domain of \cref{step:diag}, plus a health
    flag; never a premise for safety \\[3pt]
row key
  & one row per entry and coalition, over the derived closure
    (\cref{step:coalitions}) \\[3pt]
witness
  & replayable, and canonical enough that the discovery route is not
    observable in it (\cref{step:verdict}) \\
\bottomrule
\end{tabular}
\end{center}

The three-status split is the part most likely to be collapsed under time
pressure, and \Conditional{} is precisely the projection of ``model proved,
deployment open.'' It is a derived value rather than a primitive alongside the
other three statuses --- which is why
an obligation naming a target fact and one naming a backend fact should not share
a field.

\subsection{References}
\label{sec:refs}

Verified against Crossref, dblp, and author- or proceedings-hosted copies.
Publisher pages for USENIX, IEEE, ACM, and Springer were not directly
retrievable, so those entries rest on the registration authority rather than the
publisher's own page.

\begingroup
\small
\begin{description}[leftmargin=1.2em,style=unboxed,itemsep=2pt]

\item[Self-composition and 2-safety.]
G.~Barthe, P.~R. D'Argenio, T.~Rezk. \emph{Secure Information Flow by
Self-Composition}. CSFW 2004, 100--114; extended in \emph{Math.\ Struct.\ in
Comp.\ Sci.}\ 21(6):1207--1252, 2011.\quad
T.~Terauchi, A.~Aiken. \emph{Secure Information Flow as a Safety Problem}. SAS
2005, 352--367.\quad
M.~R. Clarkson, F.~B. Schneider. \emph{Hyperproperties}. CSF 2008, 51--65;
\emph{J.\ Computer Security} 18(6):1157--1210, 2010.\quad
N.~Benton. \emph{Simple Relational Correctness Proofs for Static Analyses and
Program Transformations}. POPL 2004, 14--25.

\item[Information-flow type systems.]
D.~M. Volpano, C.~E. Irvine, G.~Smith. \emph{A Sound Type System for Secure Flow
Analysis}. \emph{J.\ Computer Security} 4(2--3):167--187, 1996.\quad
A.~Sabelfeld, A.~C. Myers. \emph{Language-Based Information-Flow Security}.
\emph{IEEE JSAC} 21(1):5--19, 2003.

\item[Constant-time verification.]
J.~B. Almeida, M.~Barbosa, G.~Barthe, F.~Dupressoir, M.~Emmi. \emph{Verifying
Constant-Time Implementations}. USENIX Security 2016, 53--70. \emph{(ct-verif;
toolchain unmaintained since 2017, benchmark corpus still current.)}\quad
Z.~Zhang, G.~Barthe. \emph{CT-LLVM: Automatic Large-Scale Constant-Time
Analysis}. IACR ePrint 2025/338.\quad
L.-A. Daniel, S.~Bardin, T.~Rezk. \emph{Binsec/Rel}. IEEE S\&P 2020, 1021--1038;
extended in \emph{ACM TOPS} 26(2), 2023.\quad
L.~Cai, F.~Song, T.~Chen. \emph{Towards Efficient Verification of Constant-Time
Cryptographic Implementations}. FSE 2024.

\item[Compiler preservation of hyperproperties.]
G.~Barthe, S.~Blazy, B.~Gr\'egoire, R.~Hutin, V.~Laporte, D.~Pichardie,
A.~Trieu. \emph{Formal Verification of a Constant-Time Preserving C Compiler}.
\emph{PACMPL} 4(POPL), Art.~7, 2020.\quad
J.~Rosemann, S.~Hack, D.~Garg. \emph{Non-interference Preserving Optimising
Compilation}. \emph{PACMPL} 9(OOPSLA2):1429--1454, 2025.\quad
C.~Abate, R.~Blanco, D.~Garg, C.~Hri\c{t}cu, M.~Patrignani, J.~Thibault.
\emph{Journey Beyond Full Abstraction}. CSF 2019, 256--271.\quad
M.~Patrignani, A.~Ahmed, D.~Clarke. \emph{Formal Approaches to Secure
Compilation}. \emph{ACM Computing Surveys} 51(6), 2019.

\item[Mechanized LLVM and MLIR.]
J.~Zhao, S.~Nagarakatte, M.~M.~K. Martin, S.~Zdancewic. \emph{Formalizing the
LLVM Intermediate Representation for Verified Program Transformations}. POPL
2012, 427--440.\quad
Y.~Zakowski, C.~Beck, I.~Yoon, I.~Zaichuk, V.~Zaliva, S.~Zdancewic.
\emph{Modular, Compositional, and Executable Formal Semantics for LLVM IR}.
\emph{PACMPL} 5(ICFP), 2021. \emph{(current Vellvm.)}\quad
L.~Silver, P.~He, E.~Cecchetti, A.~K. Hirsch, S.~Zdancewic. \emph{Semantics for
Noninterference with Interaction Trees}. ECOOP 2023, 29:1--29:29.\quad
N.~P. Lopes, J.~Lee, C.-K. Hur, Z.~Liu, J.~Regehr. \emph{Alive2: Bounded
Translation Validation for LLVM}. PLDI 2021, 65--79.\quad
M.~Fehr, Y.~Fan, H.~Pompougnac, J.~Regehr, T.~Grosser. \emph{First-Class
Verification Dialects for MLIR}. \emph{PACMPL} 9(PLDI):1466--1490, 2025.\quad
S.~Bhat, A.~Keizer, C.~Hughes, A.~Goens, T.~Grosser. \emph{Verifying Peephole
Rewriting in SSA Compiler IRs}. ITP 2024, 9:1--9:20. \emph{(lean-mlir.)}

\item[Constant-time by construction.]
J.~B. Almeida et al. \emph{Jasmin: High-Assurance and High-Speed Cryptography}.
CCS 2017, 1807--1823; \emph{The Last Mile}, IEEE S\&P 2020, 965--982.\quad
S.~Cauligi et al. \emph{FaCT: A DSL for Timing-Sensitive Computation}. PLDI
2019, 174--189.

\end{description}
\endgroup

\subsection{Reading the arc backwards}

One last framing, for a reader who wants the shortest possible summary of
Part~\ref{part:walkthrough}.

The checker never proves a program safe. It reduces ``is this program safe'' to
``is this product unreachable,'' having first established that the question is
well-posed: the policy is total, the artifact is in the fragment and is the one
that shipped, and the admitted domain is non-empty. Every one of those
preconditions is a place where a plausible shortcut yields a confident wrong
answer, and the corpus exists because each shortcut is more tempting than it
looks.
